\section{Eleven integers at degree 464}\label{sec:eleven}

This branch minimizes a different part of the complexity vector. It
retains eleven coordinates and the complete native source34 coefficient
tuple, but combines two incoming constructions in a way that substantially
reduces degree. The three-root coprime projection of \cite{coprime}
supplies the variable count. Powered saturation and supported-modulus
interpolation from \cite{round14} allow a much smaller scalar carrier.
Two elementary changes to the guard and the radical weights then reduce
the literal arithmetic cost.

\subsection{The complete equation}

The eleven independent integers are
\[
 j,C,z,h,r,v,w,T,c,n,s.
\]
The six fixed coefficients are denoted by
$\beta,K_0,M_C,M_F,A_0,B_0$. They are the unchanged source34 tuple,
with
\[
 \beta=2^{1064921}-1,\qquad A_0=2^{1113},\qquad B_0=2^{666104}.
\]
The complete bytes of $K_0,M_C,M_F$, as well as the other three
coefficients, are supplied with their bindings and finite reconstruction
recipe. In particular, $M_F$ already includes the source shift. None of
these six coefficients is a free parameter or an existential unknown.

Define the following polynomial abbreviations:
\begin{align}
 J&=6j+1,\qquad q=\beta J+1,\qquad W=12w+4,\nonumber\\
 X&=qW,\qquad Y=(2v+1)q^3,\qquad E=XY,\nonumber\\
 a&=E+Y,\qquad H=4a+3,\nonumber\\
 \Delta&=a^2+H=(a+2)^2-1,\label{eq:outer-one}\\
 F&=(K_0+W)C-z(q-1),\nonumber\\
 R&=(q(q-F)-C)(q^2-1)+(M_C+qM_F)J,\nonumber\\
 k&=R+1+hE,\qquad \eta=c-Yk,\label{eq:outer-two}\\
 D_{\rm err}&=(A_0F-q)^2+(B_0C-q)^2,\nonumber\\
 \mathcal B&=k^2(q^2-D_{\rm err})-4q^2\eta^2,\qquad D=3\mathcal B-4,\label{eq:budget}\\
 L&=XY^2k,\qquad D_t=1+L(L+k),\nonumber\\
 D_g&=1+D D_t s^2,\qquad P=D_tD_g.\label{eq:guard}
\end{align}
The new strong and carrier scales, and the auxiliary data, are
\begin{align}
 I&=2rc^2, & S&=srH,\label{eq:new-scales}\\
 D_f&=1+\Delta I^2, & Q&=\Delta^2 I^2, & \delta&=Q-1,\nonumber\\
 V&=D_f(c-R+4IT)-c,\nonumber\\
 C_y&=QV^2-1, & D_a&=\delta C_y,\label{eq:auxiliary}\\
 K&=X(3c-2X)+2, & Z_c&=HT+K,\nonumber\\
 A_c&=(Z_c^2+3H^2)(Z_c^3-19H^3).\label{eq:carrier}
\end{align}
For the centered equation set
\begin{equation}
 Z=S(n+QV)-A_c,\qquad
 u=4S^2D_f,\qquad b=S^2D_a,\qquad e=16S^2P.
 \label{eq:three-roots}
\end{equation}
Finally define
\begin{equation}
 \boxed{\displaystyle
 \mathcal E_{464}=
 \bigl[(Z^2-u-b-e)^2-4(ub+ue+be)\bigr]^2
       -64ubeZ^2=0.}
 \label{eq:E464}
\end{equation}
Equations \eqref{eq:outer-one}--\eqref{eq:E464} specify a single
ordinary integer polynomial. Abbreviations do not introduce coordinates,
and all exponents are fixed nonnegative integers. The literal circuit
contains 24 additions, 21 subtractions, and 77 multiplications: 122 live
binary operations in total. The direct companion replaces $Z$ by
$Sn-A_c$ and has exact degree 488 in 121 operations.

The universal identity underlying \eqref{eq:E464} is
\begin{equation}
 \mathcal E=
 \prod_{(\epsilon_1,\epsilon_2,\epsilon_3)\in\{-1,1\}^3}
 (Z-\epsilon_1\sqrt u-\epsilon_2\sqrt b-\epsilon_3\sqrt e).
 \label{eq:three-norm}
\end{equation}
It holds algebraically even for zero, negative, or dependent radicands.
One may prove it by successive quadratic conjugation, or work in the
free algebra generated by three formal square roots. The checker
verifies its entire polynomial expansion independently.

\subsection{The inherited source interface}

The old centered proof supplies four precise arithmetic interfaces,
whose statements and full ordinary source proofs are retained in the
package. They are used here in the order stated.

First, the disk inequality $\mathcal B\ge2$ yields the common strict
sign of $q,F,C,q-F-C$ and
\begin{equation}
 k\ne0,\quad D_{\rm err}<q^2,\quad
 \left|\frac{c}{k}-Y\right|<\frac12,\quad
 \sgn(c)=\sgn(kY).
 \label{eq:centered-strip}
\end{equation}
Second, the strong and auxiliary Pell equations, together with
$c^2\mid I$, recover the missing main Pell ordinate by absolute index
folding. Third, the \emph{pre-projection} part of the signed rank
argument yields
\begin{equation}
 p=R=2t+1,\quad t\ge24,\quad
 |Y|>\frac{(|X|-3)^t}{2}-1,\quad
 |H|>|X|(|X|-3)^t.
 \label{eq:rank-bound}
\end{equation}
It does not assume the main scalar projection. Fourth, after that
projection has been recovered, the source orientation and half-binomial
arguments regenerate a valid old native certificate and hence a finite
syntactic refutation. Using the fourth interface before proving the
projection would be circular; this construction does not do so.

For completeness, a canonical source witness from a refutation supplies
$a+2>1$ even, $R\equiv3\pmod4$, $c\equiv R\pmod4$, a positive
centered budget, and the old main projection. The fixed-pair interpolation
theorems described below apply to these canonical witnesses. These
retained source facts are substantial inherited mathematics, rather
than consequences of finite polynomial testing.

\subsection{Bootstrap without a hidden nonzero assumption}

The literal coefficient tuple gives, for every integer assignment,
\[
 q\equiv0\pmod2,\quad q\equiv2\pmod3,\quad
 |q|\ge\beta-1,\quad |X|\ge4|q|,\quad |Y|\ge|q|^3.
\]
In particular $H$ is odd and nonzero, and $\Delta>1$ with
$\Delta\equiv3\pmod4$. Writing $\ell=XY^2$ gives
$D_t=1+\ell(\ell+1)k^2\ge1$. It is odd, and it is one modulo three
because $\ell$ is zero or two modulo three.

\begin{lemma}[The reduced guard]\label{lem:guard}
The budget $\mathcal B$ is divisible by four for every integer
assignment. At a zero of \eqref{eq:E464}, one has
$S\ne0$, $s r c I\ne0$, $\mathcal B\ge2$, $D_g>0$, and $D_a>0$.
\end{lemma}
\begin{proof}
The integers $q,A_0,B_0$ are even. Therefore both $q^2$ and
$D_{\rm err}$ are divisible by four, as is the complete expression
for $\mathcal B$ in \eqref{eq:budget}. It follows that $D$ is divisible
by four and $D_g\equiv1\pmod4$. This makes the old external factor
four in $1+4DD_t s^2$ unnecessary.

If $S=0$, then $u=b=e=0$, and the polynomial is $A_c^8$.
Let $\pi(x)=(x^2+3)(x^3-19)$. Since
\[
 A_c=H^5\pi(T+K/H),
\]
it cannot vanish: the quadratic factor has no real root and the cubic
has no rational root. Thus $S\ne0$, so $s,r\ne0$.

At this point $c\ne0$ must not yet be assumed. If $c=0$, then
$I=Q=0$, $D_f=1$, $\delta=C_y=-1$, and $D_a=1$.
If $c\ne0$, then $|I|\ge2c^2$ and $D_f>|c|+1$. In the integer
expression $V=D_f(c-R+4IT)-c$, either the parenthesized integer is zero,
giving $|V|=|c|\ge1$, or $|V|\ge D_f-|c|>1$. Hence $Q>1$,
$C_y>0$, and $D_a>0$. Thus $D_a$ is positive in both cases.

If $D_g<0$, its radical in \eqref{eq:three-norm} is the sole nonzero
imaginary term: $D_f,D_a$ are positive, $D_t\ge1$, and $S\ne0$.
No conjugate factor can vanish. Since $D_g\equiv1\pmod4$, it follows
that $D_g>0$. If $\mathcal B\le1$, its divisibility by four gives
$\mathcal B\le0$, so $D\le-4$ and $D_g\le-3$, a contradiction.
Consequently $\mathcal B\ge2$.

The disk argument now proves \eqref{eq:centered-strip}. Because $Y$ is
a nonzero integer, that strip excludes $c=0$. Thus $I\ne0$ as well.
\end{proof}

\subsection{Separating three radicals and recovering the main root}

At a zero, divide the vanishing conjugate factor by $S$. After a choice
of signs, the rational number on its left is a signed sum
\begin{equation}
 2\sqrt{D_f}+\sqrt{D_a}+4\sqrt{P}.
 \label{eq:weights}
\end{equation}
All three radicands are positive odd integers. Indeed $I$ is even, so
$D_f\equiv D_a\equiv1\pmod4$, while $D_t,D_g$ are odd.
For each odd integer radicand, factor its square root as an odd integer
times the square root of its odd squarefree part. In any nontrivial
squareclass occurring in \eqref{eq:weights}, the coefficients have
distinct 2-adic valuations drawn from $1,0,2$. The term of least
valuation cannot cancel with the others. Linear independence of distinct
squarefree radicals then forces all three roots to be rational, hence
integral.

The identity
\begin{equation}
 \gcd(D_t,D_g)=\gcd(D_t,1+DD_t s^2)=1
 \label{eq:coprime}
\end{equation}
shows that the positive square $P=D_tD_g$ has two separately square
factors. Thus there exist integers $f,\tau,g,\widetilde y$ with
\[
 f^2=D_f,\quad \tau^2=D_t,\quad g^2=D_g,\quad
 \widetilde y^2=D_a.
\]
Put $y=\widetilde y/\delta$ initially in $\mathbb Q$.
It satisfies $QV^2-(Q-1)y^2=1$. The retained adjacent-coefficient
Pell descent makes $y$ an integer. Its hypotheses $Q>1$, integral
$V$, and square $Q=\Delta^2I^2$ have already been established.
The signed factor now also gives $S\mid A_c$.

For clarity, the main-root recovery occurs before scalar projection.
Set $B=|a+2|$, $C_*=|c|$, and $I_*=|I|$. Strong Pell classification
writes
\[
 |f|=X_B(m_0),\qquad I_*=Y_B(m_0),\qquad |f|>2I_*.
\]
Here $X_B(j),Y_B(j)$ are the usual integer coordinates of
$(B+\sqrt{B^2-1})^j$. Classification of the auxiliary norm and its
two residues gives an odd index $\nu$ and
$Y_B(\nu)\equiv\pm C_*\pmod{|f|}$. Fold that index modulo $2m_0$
and reflect to $0\le p\le m_0$. Since
$Y_B(p)+C_*\le2I_*<|f|$, no modular wrap is possible, and
$C_*=Y_B(p)$ with $p\ge1$.

The two residues used here are unchanged:
\[
 V\equiv-c\pmod f,\qquad V\equiv-R\pmod c.
\]
The second follows directly from
$V+R=\Delta I^2(c-R)+4ID_fT$ and $c^2\mid I$.
Lucas composition supplies the required rank divisibility with an even
multiplier $I/c^2$; no odd-multiplier normalization is assumed. The
retained pre-projection argument now applies and proves
\eqref{eq:rank-bound}.

\subsection{Powered divisibility replaces a large carrier}

The old construction used $S=IH^5$ to force $H\mid K$. Our smaller
$S=srH$ forces only $H\mid K^5$. The exact coefficient-size bounds
are strong enough to recover the same final projection from this weaker
statement.

Let
\[
 K_{\rm quad}=X(X+2ac)-1,\qquad \sigma=\sgn(c),\qquad
 \Pi=(\sigma X-2^R)(2^R\sigma X+1).
\]
The literal identity
\begin{equation}
 K=HXc-2K_{\rm quad}
 \label{eq:short-projector}
\end{equation}
is verified by expansion. Since $H\mid S\mid A_c$ and
$A_c\equiv K^5\pmod H$, we obtain $H\mid K^5$. Oddness of $H$
then gives $H\mid K_{\rm quad}^5$. The signed main Pell projection
identity supplied by the recovered main ordinate is
\[
 2^R K_{\rm quad}\equiv\Pi\pmod H.
\]
Raising this congruence to the fifth power gives $H\mid\Pi^5$.
Only powers of two, which are units modulo the odd $H$, are cancelled.

\begin{lemma}[The fifth power is smaller than the modulus]
Under the established pre-projection hypotheses, $|\Pi|^5<|H|$.
\end{lemma}
\begin{proof}
Put $x=|X|$. The unchanged actual coefficient tuple gives $x\ge2^{64}$,
far more than is needed. With $R=2t+1$ and $t\ge24$, elementary
estimates give
\[
 |\Pi|^5<2^{20t+20}x^{10},\qquad
 |H|>x(x-3)^t\ge\frac{x^{t+1}}{2^t}.
\]
The first inequality follows by bounding each of the two factors in
$\Pi$ by $2^{R+1}x$. The comparison needed for the two displayed
bounds is $x^{t-9}>2^{21t+20}$. It follows from
\[
 64(t-9)-(21t+20)=43t-596\ge436>0.
\]
\end{proof}

Divisibility and the strict height bound force $\Pi=0$. The second
factor of $\Pi$ is odd and therefore nonzero, so
$\sigma X=2^R$. The retained orientation and half-binomial proof can
now be used: it gives $X,c,k,Y>0$, the required strict native strip,
and a fresh valid old certificate. This completes soundness of the new
equation relative to the inherited native calibration.

\subsection{Completeness with a fixed strong pair}\label{sec:interpolation}

Begin with a canonical outer witness supplied by a finite syntactic
refutation in the retained source theorem. Its $\mathcal B$ is positive,
$D_t=\tau^2$, and $DD_t$ is a positive nonsquare: it is two modulo
three. Ordinary Pell existence supplies positive $g,s$ such that
\[
 g^2-DD_t s^2=1.
\]
This directly gives the reduced guard $D_g=g^2$. There is no need to
square this unit to create an extra factor four.

Choose one positive strong pair $I,f$ with $4sc^2\mid I$ and put
$r=I/(2c^2)$. Then $r>0$ and $2s\mid r$. Every prime divisor of
$M=sr$ divides $I$. Hold this strong pair, $r$, and $M$ fixed.

The supported-modulus interpolation theorem of \cite{round14} says
that every residue class of $T$ modulo any such $M$ is realized by an
integer auxiliary Pell witness $y$ for this fixed strong pair. It
applies to the same canonical even base and the same auxiliary chart
used here. Its proof does not depend on the choice of determinant or
centered-budget guard. The independent review checks the hypotheses
$a+2$ even, $R\equiv c\equiv3\pmod4$, and $4sc^2\mid I$ explicitly.

We recall the mechanism because the order of quantifiers matters.
Writing $I=Y_A(m_0)$ and $f=X_A(m_0)$, the exact period needed modulo
$f^2$ is $\lambda=4m_0f$. Set $d=\operatorname{lcm}(\lambda,4I)$ and
vary a starting auxiliary index through $\nu(h)=\nu_0+dh$.
The auxiliary coordinate
\begin{equation}
 T(h)=\frac{X_{\Delta I}(\nu(h))/(\Delta I)+c-f^2(c-R)}{4If^2}
 \label{eq:interp-coordinate}
\end{equation}
is an integer. Its denominator is part of the interpolation proof, not
part of the final polynomial. For every prime $p\mid I$, Pell addition
and valuation lifting give
\begin{equation}
 v_p\bigl(T(h_1)-T(h_2)\bigr)=v_p(h_1-h_2).
 \label{eq:isometry}
\end{equation}
The exact rank valuations give $v_p(d)=v_p(4I)$; the factor $f$ in
$\lambda$ cannot be omitted. For $p=2$, the first-coordinate difference
has valuation $2(v_2(\Delta I)+v_2(d(h_1-h_2)))-1$, strictly larger
than the ordinate contribution because the index difference is
divisible by four. Thus the two-adic case is included separately.
Equation \eqref{eq:isometry} induces a permutation of each supported
prime-power residue ring. Chinese remaindering of the corresponding
$h$ classes proves the full supported-modulus statement.

The polynomial $\pi(x)=(x^2+3)(x^3-19)$ has a root modulo every positive
integer. The inherited elementary proof treats the exceptional primes
$2,3,19$, uses lifting at the other primes, and combines the roots by
CRT. Since the canonical projection gives $H\mid K$, write
$\kappa=K/H$. Choose a residue with
\[
 \pi(T+\kappa)\equiv0\pmod M,
\]
and only then choose its interpolated auxiliary witness. It follows that
\[
 \frac{A_c}{S}=
 \frac{H^4\pi(T+\kappa)}{sr}\in\mathbb Z.
\]
Finally set
\begin{equation}
 n=\frac{A_c}{S}+2f+\delta y+4\tau g-QV.
 \label{eq:choose-n}
\end{equation}
One conjugate factor in \eqref{eq:three-norm} now vanishes, giving an
integer zero of $\mathcal E_{464}$.

The witness order is outer native data, guard Pell unit, one strong
pair, carrier residue, auxiliary $T,y$, and finally $n$. The strong
pair is never changed after the residue is prescribed. The new
construction preserves existence of zeros; it does not identify
old and new auxiliary tuples pointwise.

\subsection{Exact degree and paid arithmetic cost}

The principal degree ledger is
\begin{center}
\begin{tabular}{lrrrrrrrr}
\toprule
Expression &$I$&$S$&$D_f$&$Q$&$V$&$D_a$&$P$&$A_c$\\
Degree &3&8&18&30&23&106&94&35\\
\midrule
Expression &$u$&$b$&$e$&$Z_0=SQV$&$Z_1=Sn-A_c$\\
Degree &34&122&110&61&35\\
\bottomrule
\end{tabular}
\end{center}
The center uses the already computed $QV$. Its exact cancellation is
\begin{equation}
 b-Z_0^2=S^2(-QV^2-Q+1).
 \label{eq:center-cancel}
\end{equation}
Thus $\omega=(Z_0+Z_1)^2-b$ has degree at most 96. Rewriting the
identical polynomial as
\begin{equation}
 \bigl[(\omega-u-e)^2-4(ue+b(u+e))\bigr]^2
       -64ube(Z_0+Z_1)^2
 \label{eq:reduced-degree}
\end{equation}
shows that the unique highest-degree part inside the outer square is
$-4b_{122}e_{110}$, of degree 232. The last product has degree
$34+122+110+122=388$. Therefore the degree is exactly 464, with leading
form $16b_{122}^2e_{110}^2$.

The independently checked explicit coefficient
\begin{equation}
 [j^{260}w^{56}v^{64}C^4s^{12}r^{24}c^{32}h^{12}]\mathcal E_{464}
   =2^{484}3^{318}\beta^{260}>0
 \label{eq:leading464}
\end{equation}
proves attainment for the actual fixed tuple. The exponent sum is 464.
For the direct companion, $Z$ has degree 35 and the leading part is
$b_{122}^4$, of degree 488. An attained coefficient is
\[
 [j^{248}w^{64}v^{56}C^8s^8r^{40}c^{64}]\mathcal E_{488}
   =2^{480}3^{312}\beta^{248}>0.
\]
The degree certificate uses exact leading homogeneous arithmetic and
the cancellation identity, not a numerical evaluation of the polynomial
at a few points.

The delivered graphs count every live arithmetic node, authenticate all
six coefficient payloads, and reconstruct the coefficient recipe.
The 122-operation count is separate from the cancellation-aware degree
argument: an uncollected syntactic degree bound for the shortest graph
would overestimate the true degree.

\begin{theorem}[The eleven-integer equation]\label{thm:eleven}
For the six fixed coefficients supplied with the packet,
$\mathcal E_{464}$ has eleven unrestricted integer unknowns, exact
degree 464, and a 122-operation division-free evaluation circuit.
The retained source calibration and the arguments above prove
\[
 \ZFC\vdash
 [\forall x\in\mathbb Z^{11}\;\mathcal E_{464}(x)\ne0]
 \longleftrightarrow\Con(\ZFC).
\]
The same conclusion holds for the degree-488, 121-operation direct
companion.
\end{theorem}
